\section{Geometric and Combinatorial Preliminaries}
\label{sec:preliminaries}

We begin by introducing the geometric foundation of the simplex memory architecture. The construction proceeds in two steps: first we define the continuous regular simplex, then we discretize it by a lattice of points that will serve as the hidden nodes of the network.

\begin{definition}[The regular $n$-simplex]
For $n\geq 2$, the standard regular $n$-simplex is the subset of $\R^{n+1}$ defined by
\[
\Delta^n
:=
\left\{x=(x_0,\dots,x_n)\in \R^{n+1}: x_i\geq 0,\ \sum_{i=0}^n x_i=1\right\}.
\]
This is an $n$-dimensional polytope lying in the affine hyperplane $\sum_{i=0}^n x_i=1$. Its vertices are the canonical basis vectors
\[
e_0, e_1, \dots, e_n \in \R^{n+1}.
\]
Throughout this work we distinguish two vertices and relabel them as
\[
a := e_0,
\qquad
b := e_1,
\qquad
c_1 := e_2,\ \dots,\ c_{n-1}:=e_n.
\]
Thus $a$ and $b$ play the roles of the two extreme vertices (the ``backbone'' endpoints), while $c_1,\dots,c_{n-1}$ form the intermediate vertex set.
\end{definition}

\begin{figure}[ht]
\centering
\begin{tikzpicture}[scale=1.0, >=stealth, line width=0.5pt,
                    line cap=round, line join=round]
  % Oblique projection: x_0 at 225deg, x_1 at 0deg, x_2 at 90deg
  % Vertices at projected distance 1.5; all three axes equal projected length 2.8
  % x_0 tip: 2.8*(cos225, sin225) = (-1.980, -1.980)
  \coordinate (O) at (0, 0);
  \coordinate (A) at (-0.55, -0.55);  % a   = e_0  (on 225deg axis, proj dist ~0.78)
  \coordinate (B) at ( 1.5,   0   );  % b   = e_1
  \coordinate (C) at ( 0,     1.5 );  % c_1 = e_2

  % 1. Triangle fill (drawn first so dashed lines sit on top)
  \fill[gray!10] (A) -- (B) -- (C) -- cycle;

  % 2. Hidden axis segments (O to each vertex, passes through face): dashed
  \draw[dashed] (O) -- (A);
  \draw[dashed] (O) -- (B);
  \draw[dashed] (O) -- (C);

  % 3. Visible axis segments (beyond each vertex to tip): solid with arrow
  \draw[->] (A) -- (-1.40, -1.40) node[below left]  {$x_0$};
  \draw[->] (B) -- ( 2.8,   0     ) node[right]        {$x_1$};
  \draw[->] (C) -- ( 0,     2.8   ) node[above]        {$x_2$};

  % 4. Triangle edges
  \draw (A) -- (B) -- (C) -- cycle;

  % Vertex labels
  \node[below left=2pt]  at (A) {$a$};
  \node[below right=2pt] at (B) {$b$};
  \node[above right=2pt] at (C) {$c_1$};

  % Facet labels
  \node[left=4pt]  at ($(A)!0.5!(C)$) {$F_{\mathrm{in}}$};
  \node[right=4pt] at ($(B)!0.5!(C)$) {$F_{\mathrm{out}}$};
  \node[above left=1pt, font=\scriptsize] at (C) {$F_{\mathrm{mid}}$};
\end{tikzpicture}
\caption{The regular $2$-simplex $\Delta^2 \subset \mathbb{R}^3$, with vertices
  $a=e_0$, $b=e_1$, $c_1=e_2$ on the coordinate axes.
  The input facet is $F_{\mathrm{in}}=\operatorname{conv}(a,c_1)$,
  the output facet $F_{\mathrm{out}}=\operatorname{conv}(c_1,b)$,
  and their intersection $F_{\mathrm{mid}}=\{c_1\}$.}
\label{fig:simplex-geometry}
\end{figure}

The geometric structure becomes a neural architecture only after discretization. We now introduce the lattice that provides the discrete set of hidden nodes.

\begin{definition}[The discrete simplex lattice]
Fix an integer $m\geq 2$ and write $s := m-1$. The lattice of hidden nodes is defined by
\[
V_{n,m}
:=
\left\{\alpha=(\alpha_0,\dots,\alpha_n)\in \Z_{\geq 0}^{n+1}: \sum_{i=0}^n \alpha_i = s\right\}.
\]
Each $\alpha\in V_{n,m}$ corresponds to the geometric point
\[
x(\alpha)=\frac{1}{s}\alpha \in \Delta^n.
\]
Equivalently, $V_{n,m}$ is the set of discrete barycentric coordinates of the simplex lattice with exactly $m$ points (including endpoints) on every one-dimensional edge.
\end{definition}

The input and output interfaces of the architecture are realized by two opposite facets of the simplex. Their intersection forms a shared intermediate buffer.

\begin{definition}[Input, output, and shared facets]
We define the input facet, output facet, and shared intermediate face by
\[
F_{\mathrm{in}} := \{\alpha\in V_{n,m}: \alpha_1=0\},
\qquad
F_{\mathrm{out}} := \{\alpha\in V_{n,m}: \alpha_0=0\},
\]
\[
F_{\mathrm{mid}} := F_{\mathrm{in}}\cap F_{\mathrm{out}} = \{\alpha\in V_{n,m}: \alpha_0=\alpha_1=0\}.
\]
Geometrically, these correspond to the facets
\[
\conv(a,c_1,\dots,c_{n-1}),
\qquad
\conv(c_1,\dots,c_{n-1},b),
\]
and their common face
\[
\conv(c_1,\dots,c_{n-1}).
\]
\end{definition}

We now compute the cardinalities of these sets. These quantities determine the number of hidden nodes, the size of the input and output interfaces, and the scaling of the architecture with respect to the parameters $(n,m)$.

\begin{proposition}[Basic cardinalities]
The lattice sizes are given by
\[
\card{V_{n,m}} = \binom{m+n-1}{n},
\qquad
\card{F_{\mathrm{in}}} = \card{F_{\mathrm{out}}} = \binom{m+n-2}{n-1},
\]
\[
\card{F_{\mathrm{mid}}} = \binom{m+n-3}{n-2}.
\]
\end{proposition}

\begin{proof}
All three formulas follow from standard stars-and-bars counting. The total number of hidden nodes equals the number of nonnegative integer solutions of
\[
\alpha_0+\cdots+\alpha_n = m-1,
\]
which is $\binom{m+n-1}{n}$. For the input facet, we impose $\alpha_1=0$, leaving
\[
\alpha_0+\alpha_2+\cdots+\alpha_n=m-1,
\]
with $n$ variables; hence $\card{F_{\mathrm{in}}}=\binom{m+n-2}{n-1}$. The output facet is identical by symmetry. Finally, the shared face is obtained by imposing both $\alpha_0=0$ and $\alpha_1=0$, leaving $n-1$ variables and thus
\[
\card{F_{\mathrm{mid}}}=\binom{m+n-3}{n-2}.
\]
\end{proof}

\begin{remark}[Interface-to-volume ratio]
The proportion of lattice nodes lying on the input facet (and likewise on the output facet) is
\[
\frac{\card{F_{\mathrm{in}}}}{\card{V_{n,m}}}
=
\frac{\binom{m+n-2}{n-1}}{\binom{m+n-1}{n}}
=
\frac{n}{n+m-1}.
\]
This ratio suggests a structural interpretation. Increasing $m$ deepens the simplex lattice while reducing the facet proportion, whereas increasing $n$ enlarges the relative size of the input and output interfaces. In this sense, the pair $(n,m)$ naturally controls the balance between broad interfaces and deep interior propagation.
\end{remark}
